\chapter{2001年京蒙皖卷（春文）}
\section{单选题}
\begin{problem}
集合 \(M = \{1, 2, 3, 4, 5\}\) 的子集个数是（\quad）

\choices
    {\(32\)}
    {\(31\)}
    {\(16\)}
    {\(15\)}
\end{problem}

\begin{answer}
A
\end{answer}

\begin{solution}
集合 \(M\) 有 \(5\) 个元素. 对于 \(M\) 中的每一个元素，选入子集或不选入子集各有 \(2\) 种选择，且这些选择彼此独立，因此子集总数为
\[
2^5=32
\]
所以选 A.
\end{solution}

\begin{problem}
函数 \(f(x) = a^x\)（\(a > 0\)，\(a \neq 1\)）对于任意的实数 \(x\)，\(y\) 都有（\quad）

\choices
    {\(f(xy) = f(x)f(y)\)}
    {\(f(xy) = f(x) + f(y)\)}
    {\(f(x + y) = f(x)f(y)\)}
    {\(f(x + y) = f(x) + f(y)\)}
\end{problem}

\begin{answer}
C
\end{answer}

\begin{solution}
由指数运算法则，\(f(x+y)=a^{x+y}=a^x a^y=f(x)f(y)\)，所以第三个选项正确. 例如取 \(x=y=0\)，第二个和第四个选项都变成 \(1=2\)，不成立；而第一个选项要求对任意 \(x\)，\(y\) 有 \(a^{xy}=a^{x+y}\)，取 \(x=2\)，\(y=3\) 得 \(a^6=a^5\)，这与 \(a\neq1\) 矛盾. 因此正确选项为第三个选项.
\end{solution}

\begin{problem}
\(\displaystyle\lim_{n\to\infty}\frac{\mathrm C_{2n}^{n}}{\mathrm C_{2n+2}^{n+1}}=\)（\quad）

\choices
    {\(0\)}
    {\(2\)}
    {\(\frac{1}{2}\)}
    {\(\frac{1}{4}\)}
\end{problem}

\begin{answer}
D
\end{answer}

\begin{solution}
\[
\frac{\mathrm C_{2n}^{n}}{\mathrm C_{2n+2}^{n+1}}
=\frac{\dfrac{(2n)!}{(n!)^2}}{\dfrac{(2n+2)!}{((n+1)!)^2}}
=\frac{(n+1)^2}{(2n+2)(2n+1)}
=\frac{n+1}{2(2n+1)}
\]
令 \(n\to\infty\)，得
\[
\lim_{n\to\infty}\frac{n+1}{2(2n+1)}=\frac14
\]
所以选 D.
\end{solution}

\begin{problem}
函数 \(y = -\sqrt{1 - x}\)（\(x \leqslant 1\)）的反函数是（\quad）

\choices
    {\(y = x^{2} - 1\)（\(-1 \leqslant x \leqslant 0\)）}
    {\(y = x^{2} - 1\)（\(0 \leqslant x \leqslant 1\)）}
    {\(y = 1 - x^{2}\)（\(x \leqslant 0\)）}
    {\(y = 1 - x^{2}\)（\(0 \leqslant x \leqslant 1\)）}
\end{problem}

\begin{answer}
C
\end{answer}

\begin{solution}
原函数满足 \(y=-\sqrt{1-x}\)，所以 \(y\leqslant0\). 交换 \(x\)，\(y\) 得
\[
x=-\sqrt{1-y}
\]
从而 \(x\leqslant0\)，并且平方后有 \(x^2=1-y\)，即 \(y=1-x^2\). 因此反函数为 \(y=1-x^2\)（\(x\leqslant0\)），选 C.
\end{solution}

\begin{problem}
已知 \(F_{1}\)，\(F_{2}\) 是椭圆 \(\frac{x^{2}}{16} + \frac{y^{2}}{9} = 1\) 的两焦点，过点 \(F_{2}\) 的直线交椭圆于点 \(A\)，\(B\)，若 \(\left|AB\right| = 5\)，则 \(\left|AF_{1}\right| + \left|BF_{1}\right| =\)（\quad）

\choices
    {\(11\)}
    {\(10\)}
    {\(9\)}
    {\(16\)}
\end{problem}

\begin{answer}
A
\end{answer}

\begin{solution}
椭圆的长半轴为 \(4\)，所以椭圆上任意一点 \(P\) 都满足
\[
|PF_1|+|PF_2|=2\times4=8
\]
点 \(F_2\) 在椭圆内部，且在弦 \(AB\) 上，所以 \(|AF_2|+|BF_2|=|AB|=5\). 因而
\[
|AF_1|+|BF_1|
=(8-|AF_2|)+(8-|BF_2|)
=16-5=11
\]
所以选 A.
\end{solution}

\begin{problem}
设动点 \(P\) 在直线 \(x = 1\) 上，\(O\) 为坐标原点. 以 \(OP\) 为直角边，点 \(O\) 为直角顶点作等腰 \(\mathrm{Rt}\triangle\mathrm{OPQ}\)，则动点 \(Q\) 的轨迹是（\quad）

\choices
    {圆}
    {两条平行直线}
    {抛物线}
    {双曲线}
\end{problem}

\begin{answer}
B
\end{answer}

\begin{solution}
设 \(P=(1,t)\)，则 \(\overrightarrow{OP}=(1,t)\). 等腰直角三角形的另一直角边与 \(OP\) 垂直且等长，所以 \(Q\) 有两种位置：
\[
Q=(-t,1)\quad\text{或}\quad Q=(t,-1)
\]
因此 \(Q\) 的纵坐标只能是 \(1\) 或 \(-1\)，其轨迹是两条平行直线 \(y=1\) 与 \(y=-1\)，选 B.
\end{solution}

\begin{problem}
已知 \(f(x^{6}) = \log_{2}x\)，那么 \(f(8)\) 等于（\quad）

\choices
    {\(\frac{4}{3}\)}
    {\(8\)}
    {\(18\)}
    {\(\frac{1}{2}\)}
\end{problem}

\begin{answer}
D
\end{answer}

\begin{solution}
由 \(f(x^6)=\log_2x\) 可知对数的自变量满足 \(x>0\). 要求 \(f(8)\)，令 \(x^6=8=2^3\)，得 \(x=2^{1/2}=\sqrt2\). 所以
\[
f(8)=\log_2\sqrt2=\frac12
\]
故选 D.
\end{solution}

\begin{problem}
若 \(A\)，\(B\) 是锐角 \(\triangle ABC\) 的两个内角，则点 \(P(\cos B - \sin A, \sin B - \cos A)\) 在（\quad）

\choices
    {第一象限}
    {第二象限}
    {第三象限}
    {第四象限}
\end{problem}

\begin{answer}
B
\end{answer}

\begin{solution}
因为 \(\triangle ABC\) 是锐角三角形，所以 \(A+B=\pi-C>\frac\pi2\). 又 \(A\)，\(B\) 均为锐角，故
\[
A>\frac\pi2-B
\]
从而 \(\sin A>\sin\left(\frac\pi2-B\right)=\cos B\)，于是点 \(P\) 的横坐标 \(\cos B-\sin A<0\). 同理，
\[
B>\frac\pi2-A
\]
给出 \(\sin B>\sin\left(\frac\pi2-A\right)=\cos A\)，所以纵坐标 \(\sin B-\cos A>0\). 因此 \(P\) 在第二象限，选 B.
\end{solution}

\begin{problem}
如果圆锥的侧面展开图是半圆，那么这个圆锥的顶角（圆锥轴截面中两条母线的夹角）是（\quad）

\choices
    {\(30^{\circ}\)}
    {\(45^{\circ}\)}
    {\(60^{\circ}\)}
    {\(90^{\circ}\)}
\end{problem}

\begin{answer}
C
\end{answer}

\begin{solution}
设圆锥的母线长为 \(l\)，底面半径为 \(r\)，轴截面中两条母线的夹角为 \(2\alpha\)，其中 \(0<\alpha<\dfrac{\pi}{2}\). 在轴截面中，
\[
r=l\sin\alpha
\]
圆锥侧面展开后是半径为 \(l\) 的扇形，其弧长等于圆锥底面周长. 因而扇形圆心角为
\[
\frac{2\pi r}{l}=2\pi\sin\alpha
\]
题设展开图是半圆，所以扇形圆心角为 \(\pi\)，从而
\(2\pi\sin\alpha=\pi\)，\(\sin\alpha=\frac12\)
由于 \(0<\alpha<\dfrac{\pi}{2}\)，得 \(\alpha=30^{\circ}\). 因此圆锥轴截面中两条母线的夹角为
\[
2\alpha=60^{\circ}
\]
所以选 C.
\end{solution}

\begin{problem}
若实数 \(a\)，\(b\) 满足 \(a + b = 2\)，则 \(3^{a} + 3^{b}\) 的最小值是（\quad）

\choices
    {\(18\)}
    {\(6\)}
    {\(2\sqrt{3}\)}
    {\(2\sqrt[4]{3}\)}
\end{problem}

\begin{answer}
B
\end{answer}

\begin{solution}
因为 \(3^a>0\)，\(3^b>0\)，由基本不等式得
\[
3^a+3^b\geqslant2\sqrt{3^a3^b}=2\sqrt{3^{a+b}}=2\sqrt{3^2}=6
\]
当且仅当 \(3^a=3^b\)，即 \(a=b\) 时等号成立；结合 \(a+b=2\)，得 \(a=b=1\). 所以最小值为 \(6\)，选 B.
\end{solution}

\begin{problem}
如图是正方体的平面展开图.

在这个正方体中，

\begin{circlelist}
\item \(BM\) 与 \(ED\) 平行；
\item \(CN\) 与 \(BE\) 是异面直线；
\item \(CN\) 与 \(BM\) 成 \(60^{\circ}\) 角；
\item \(DM\) 与 \(BN\) 垂直.
\end{circlelist}

以上四个命题中，正确命题的序号是（\quad）

\bitmapfigure[max width=5.4cm]{img/2001/jing_meng_wan_liberal_spring/q11_fig1.png}

\choices
    {\(\circled{1}\circled{2}\circled{3}\)}
    {\(\circled{2}\circled{4}\)}
    {\(\circled{3}\circled{4}\)}
    {\(\circled{2}\circled{3}\circled{4}\)}
\end{problem}

\begin{answer}
C
\end{answer}

\begin{solution}
将正方体棱长取为 \(1\)，把图中中间的正方形 \(ABCD\) 置于平面 \(z=0\)，折叠后可取
\(A=(0,0,0)\)，\(B=(1,0,0)\)，\(C=(1,1,0)\)，\(D=(0,1,0)\)
\(E=(0,0,1)\)，\(F=(1,0,1)\)，\(M=(1,1,1)\)，\(N=(0,1,1)\)
于是
\(\overrightarrow{BM}=(0,1,1)\)，\(\overrightarrow{ED}=(0,1,-1)\)
二者不平行，所以命题\(\circled{1}\)错误. 又
\[
\overrightarrow{CN}=(-1,0,1)=\overrightarrow{BE}
\]
所以 \(CN\parallel BE\)，命题\(\circled{2}\)错误.

由
\[
\frac{\overrightarrow{CN}\cdot\overrightarrow{BM}}
{|\overrightarrow{CN}|\,|\overrightarrow{BM}|}
=\frac{1}{\sqrt2\sqrt2}=\frac12
\]
得 \(CN\) 与 \(BM\) 所成的角为 \(60^{\circ}\)，命题\(\circled{3}\)正确. 最后
\(\overrightarrow{DM}=(1,0,1)\)，\(\overrightarrow{BN}=(-1,1,1)\)
且 \(\overrightarrow{DM}\cdot\overrightarrow{BN}=0\)，所以 \(DM\) 与 \(BN\) 所成的角为 \(90^{\circ}\)，命题\(\circled{4}\)正确. 因此正确命题为\(\circled{3}\circled{4}\)，选 C.
\end{solution}

\begin{problem}
根据市场调查结果，预测某种家用商品从年初开始的 \(n\) 个月内累积的需求量 \(S_{n}\)（\(\text{万件}\)）近似地满足 \(S_{n} = \frac{n}{90}(21n - n^{2} - 5)\)（\(n = 1\)，\(2\)，\(\cdots\)，\(12\)）. 按此预测，在本年度内，需求量超过 \(1.5\text{ 万件}\) 的月份是（\quad）

\choices
    {\(5\text{ 月}\)，\(6\text{ 月}\)}
    {\(6\text{ 月}\)，\(7\text{ 月}\)}
    {\(7\text{ 月}\)，\(8\text{ 月}\)}
    {\(8\text{ 月}\)，\(9\text{ 月}\)}
\end{problem}

\begin{answer}
C
\end{answer}

\begin{solution}
记第 \(n\) 个月的需求量为 \(T_n\). 由累计需求量的定义，
\(T_n=S_n-S_{n-1}\)，\((n=1,2,\ldots,12)\)
其中 \(S_0=0\). 由题设
\[
S_n=\frac{n}{90}(21n-n^2-5)=\frac{-n^3+21n^2-5n}{90}
\]
于是
\[
\begin{gathered}
T_n
=\frac{-n^3+21n^2-5n}{90}
-\frac{-(n-1)^3+21(n-1)^2-5(n-1)}{90}\\
=\frac{-3n^2+45n-27}{90}
=\frac{-n^2+15n-9}{30}.
\end{gathered}
\]
要求月需求量超过 \(1.5\) 万件，即
\[
\frac{-n^2+15n-9}{30}>1.5=\frac{45}{30}
\]
所以
\[
-n^2+15n-9>45
\iff n^2-15n+54<0
\iff 6<n<9
\]
在 \(n=1\)，\(2\)，\(\ldots\)，\(12\) 中，满足条件的月份编号为 \(n=7\)，\(8\)，即 7 月和 8 月.
所以选 C.
\end{solution}

\section{填空题}
\begin{problem}
已知球内接正方体的表面积为 \(S\)，那么球体积等于\fillinblank{}.
\end{problem}

\begin{answer}
\(\dfrac{\pi S\sqrt{2S}}{24}\)
\end{answer}

\begin{solution}
设正方体棱长为 \(a\). 由正方体表面积为 \(S\)，有
\(6a^2=S\)，\(a=\sqrt{\frac{S}{6}}\)
球内接正方体时，正方体的体对角线就是球的直径，因此球半径为
\[
R=\frac{\sqrt3 a}{2}
\]
所以球体积为
\[
V=\frac43\pi R^3
=\frac43\pi\left(\frac{\sqrt3 a}{2}\right)^3
=\frac{\pi\sqrt3}{2}a^3
\]
代入 \(a=\sqrt{S/6}\)，得
\[
V=\frac{\pi\sqrt3}{2}\left(\frac{S}{6}\right)^{3/2}
=\frac{\pi S\sqrt{2S}}{24}
\]
\end{solution}

\begin{problem}
椭圆 \(x^{2} + 4y^{2} = 4\) 长轴上一个顶点为 \(A\)，以 \(A\) 为直角顶点作一个内接于椭圆的等腰直角三角形，该三角形的面积是\fillinblank{}.
\end{problem}

\begin{answer}
\(\dfrac{16}{25}\)
\end{answer}

\begin{solution}
不妨取长轴顶点 \(A=(2,0)\). 设等腰直角三角形的两条直角边向量分别为
\(\bs{u}=r(\cos\varphi,\sin\varphi)\)，\(\bs{v}=r(-\sin\varphi,\cos\varphi)\)
其中 \(r>0\)，这样两向量等长且互相垂直. 因为 \(A+\bs{u}\)，\(A+\bs{v}\) 都在椭圆上，分别代入 \(x^2+4y^2=4\)，得到
\(r=-\frac{4\cos\varphi}{\cos^2\varphi+4\sin^2\varphi}\)，\(r=\frac{4\sin\varphi}{\sin^2\varphi+4\cos^2\varphi}\)
消去 \(r\)，并整理得
\[
(\cos\varphi+\sin\varphi)\bigl(4\cos^2\varphi-3\sin\varphi\cos\varphi+4\sin^2\varphi\bigr)=0
\]
第二个因子为
\(4\cos^2\varphi-3\sin\varphi\cos\varphi+4\sin^2\varphi
=4-3\sin\varphi\cos\varphi\). 由于
\(\sin\varphi\cos\varphi\leqslant\frac12\)，该因子不小于
\(4-\frac32=\frac52>0\)，所以 \(\sin\varphi=-\cos\varphi\). 由 \(r>0\) 取 \(\cos\varphi=-\frac1{\sqrt2}\)，于是
\[
r=\frac{4\sqrt2}{5}
\]
等腰直角三角形的面积为两直角边乘积的一半，因此
\[
S=\frac12r^2=\frac12\left(\frac{4\sqrt2}{5}\right)^2=\frac{16}{25}
\]
\end{solution}

\begin{problem}
已知 \(\sin^{2}\alpha + \sin^{2}\beta + \sin^{2}\gamma = 1\)（\(\alpha\)，\(\beta\)，\(\gamma\) 均为锐角），那么 \(\cos \alpha \cos \beta \cos \gamma\) 的最大值等于\fillinblank{}.
\end{problem}

\begin{answer}
\(\dfrac{2\sqrt6}{9}\)
\end{answer}

\begin{solution}
令 \(u=\sin^2\alpha\)，\(v=\sin^2\beta\)，\(w=\sin^2\gamma\)，则 \(u\)，\(v\)，\(w>0\) 且 \(u+v+w=1\). 因为三个角都是锐角，\(\cos\alpha\)，\(\cos\beta\)，\(\cos\gamma\) 均为正，所以
\[
(\cos\alpha\cos\beta\cos\gamma)^2=(1-u)(1-v)(1-w)
\]
又
\[
(1-u)+(1-v)+(1-w)=2
\]
由基本不等式得
\[
(1-u)(1-v)(1-w)\leqslant\left(\frac23\right)^3=\frac8{27}
\]
等号当且仅当 \(1-u=1-v=1-w=\frac23\)，即 \(u=v=w=\frac13\). 这时可取三个锐角使 \(\sin^2\alpha=\sin^2\beta=\sin^2\gamma=\frac13\)，条件确实能够达到. 因此
\[
\cos\alpha\cos\beta\cos\gamma\leqslant\sqrt{\frac8{27}}=\frac{2\sqrt6}{9}
\]
最大值为 \(\dfrac{2\sqrt6}{9}\).
\end{solution}

\begin{problem}
已知 \(m\)，\(n\) 是直线，\(\alpha\)，\(\beta\)，\(\gamma\) 是平面，给出下列命题：

\begin{circlelist}
\item 若 \(\alpha \perp \beta\)，\(\alpha \cap \beta = m\)，\(n \perp m\)，则 \(n \perp \alpha\) 或 \(n \perp \beta\)；
\item 若 \(\alpha \parallel \beta\)，\(\alpha \cap \gamma = m\)，\(\beta \cap \gamma = n\)，则 \(m \parallel n\)；
\item 若 \(m\) 不垂直于 \(\alpha\)，则 \(m\) 不可能垂直于 \(\alpha\) 内的无数条直线；
\item 若 \(\alpha \cap \beta = m\)，\(n \parallel m\)，且 \(n \not\subset \alpha\)，\(n \not\subset \beta\)，则 \(n \parallel \alpha\) 且 \(n \parallel \beta\).
\end{circlelist}

其中正确的命题的序号是\fillinblank{}.
\end{problem}

\begin{answer}
\(\circled{2}\)，\(\circled{4}\)
\end{answer}

\begin{solution}
逐一判断四个命题.

对于命题\(\circled{1}\)，取 \(\alpha\) 为 \(xy\) 平面，\(\beta\) 为 \(xz\) 平面，则 \(\alpha\perp\beta\)，且交线 \(m\) 为 \(x\) 轴. 取方向向量为 \((0,1,1)\) 的直线 \(n\)，则 \(n\perp m\)，但 \(n\) 的方向既不垂直于 \(\alpha\)（不平行于 \((0,0,1)\)），也不垂直于 \(\beta\)（不平行于 \((0,1,0)\)）. 所以命题\(\circled{1}\)错误.

命题\(\circled{2}\)中，\(m\) 与 \(n\) 都在平面 \(\gamma\) 内. 若 \(m\) 与 \(n\) 相交，则交点同时属于平行的两个平面 \(\alpha\)、\(\beta\)，这与 \(\alpha\parallel\beta\) 矛盾；而 \(m\) 与 \(n\) 不可能重合，因为重合直线将同时属于 \(\alpha\) 与 \(\beta\). 因此 \(m\)，\(n\) 是同一平面内的两条不相交直线，故 \(m\parallel n\)，命题\(\circled{2}\)正确.

命题\(\circled{3}\)错误. 例如取 \(\alpha\) 为 \(xy\) 平面，取 \(m\) 为其中的 \(x\) 轴. 显然 \(m\) 不垂直于 \(\alpha\). 对于任意实数 \(t\)，平面 \(\alpha\) 内过点 \((t,0,0)\) 且平行于 \(y\) 轴的直线都与 \(m\) 垂直，这样的直线有无数条.

对于命题\(\circled{4}\)，\(m=\alpha\cap\beta\)，且 \(n\parallel m\)，所以 \(n\) 的方向平行于平面 \(\alpha\) 和 \(\beta\). 若 \(n\) 与 \(\alpha\) 相交，则因 \(n\) 的方向平行于 \(\alpha\)，必有 \(n\subset\alpha\)，这与题设矛盾；故 \(n\) 与 \(\alpha\) 不相交，从而 \(n\parallel\alpha\). 同理 \(n\parallel\beta\). 命题\(\circled{4}\)正确.

因此正确的命题为\(\circled{2}\)、\(\circled{4}\).
\end{solution}

\section{解答题}
\begin{problem}
方程 \(2x^{2} + mx + n = 0\) 有实根，且 \(2\)，\(m\)，\(n\) 为等差数列的前三项. 求该等差数列公差 \(d\) 的取值范围.
\end{problem}

\begin{solution}
因为 \(2\)，\(m\)，\(n\) 是等差数列的前三项，所以
\[
m-2=n-m
\]
从而 \(n=2m-2\). 方程 \(2x^2+mx+n=0\) 有实根的充要条件是判别式不小于零，即
\[
\Delta=m^2-8n=m^2-8(2m-2)=m^2-16m+16\geqslant0
\]
解这个不等式：
\[
m\leqslant8-4\sqrt3\quad\text{或}\quad m\geqslant8+4\sqrt3
\]
又公差 \(d=m-2\)，所以
\[
d\leqslant6-4\sqrt3\quad\text{或}\quad d\geqslant6+4\sqrt3
\]
故公差 \(d\) 的取值范围为
\[
\left(-\infty,6-4\sqrt3\right]\cup\left[6+4\sqrt3,+\infty\right)
\]
\end{solution}

\begin{problem}
设函数 \(f(x) = \frac{x + a}{x + b}\)（\(a > b > 0\)），求 \(f(x)\) 的单调区间，并证明 \(f(x)\) 在其单调区间上的单调性.
\end{problem}

\begin{solution}
函数定义域为 \(\R\smallsetminus\{-b\}\)，分成区间 \((-\infty,-b)\) 和 \((-b,+\infty)\). 任取定义域内的 \(x_1<x_2\)，有
\[
f(x_1)-f(x_2)
=\frac{(x_1+a)(x_2+b)-(x_2+a)(x_1+b)}{(x_1+b)(x_2+b)}
=\frac{(a-b)(x_2-x_1)}{(x_1+b)(x_2+b)}
\]
因为 \(a>b\)，所以 \(a-b>0\). 当 \(x_1<x_2<-b\) 时，\(x_1+b\) 与 \(x_2+b\) 都为负数，分母乘积为正，故 \(f(x_1)-f(x_2)>0\). 当 \(-b<x_1<x_2\) 时，分母乘积同样为正，也有 \(f(x_1)-f(x_2)>0\). 因此 \(f(x)\) 在 \((-\infty,-b)\) 和 \((-b,+\infty)\) 上都严格单调递减.
\end{solution}

\begin{problem}
已知 \(z^{7} = 1\)（\(z \in \mathbf{C}\)，\(z \neq 1\)）.

\begin{enumerate}
\item 证明 \(1 + z + z^{2} + z^{3} + z^{4} + z^{5} + z^{6} = 0\)；
\item 设 \(z\) 的辐角为 \(\alpha\)，求 \(\cos \alpha + \cos 2\alpha + \cos 4\alpha\) 的值.
\end{enumerate}
\end{problem}

\begin{solution}
\begin{enumerate}
\item 由等比数列求和公式，
\[
1+z+z^2+z^3+z^4+z^5+z^6=\frac{z^7-1}{z-1}
\]
已知 \(z^7=1\)，且 \(z\neq1\)，所以分子为零而分母不为零，从而
\[
1+z+z^2+z^3+z^4+z^5+z^6=0
\]

\item 由 \(z^7=1\) 得 \(|z|^7=1\)，所以 \(|z|=1\). 设 \(z\) 的辐角为 \(\alpha\)，则 \(z=\cos\alpha+\i\sin\alpha\)，且 \(7\alpha=2k\pi\)，其中 \(k\in\Z\). 因此对 \(j=1\)，\(2\)，\(3\)，有
\[
\cos\left((7-j)\alpha\right)=\cos\left(2k\pi-j\alpha\right)=\cos(j\alpha)
\]
于是 \(\cos6\alpha=\cos\alpha\)，\(\cos5\alpha=\cos2\alpha\)，\(\cos4\alpha=\cos3\alpha\).
对第一问的等式取实部，得到
\[
1+\cos\alpha+\cos2\alpha+\cos3\alpha+\cos4\alpha+\cos5\alpha+\cos6\alpha=0
\]
代入上述三组相等关系，得
\[
1+2\left(\cos\alpha+\cos2\alpha+\cos4\alpha\right)=0
\]
所以
\[
\cos\alpha+\cos2\alpha+\cos4\alpha=-\frac12
\]
\end{enumerate}
\end{solution}

\begin{problem}
已知 \(VC\) 是 \(\triangle ABC\) 所在平面的一条斜线，点 \(N\) 是 \(V\) 在平面 \(ABC\) 上的射影，且在 \(\triangle ABC\) 的高 \(CD\) 上. \(AB = a\)，\(VC\) 与 \(AB\) 之间的距离为 \(h\)，点 \(M \in VC\).

\begin{enumerate}
\item 证明 \(\angle MDC\) 是二面角 \(M - AB - C\) 的平面角；
\item 当 \(\angle MDC = \angle CVN\) 时，证明 \(VC \perp\) 平面 \(AMB\)；
\item 若 \(\angle MDC = \angle CVN = \theta\)（\(0 < \theta < \frac{\pi}{2}\)），求四面体 \(MABC\) 的体积.
\end{enumerate}

\bitmapfigure[max width=0.36\linewidth]{img/2001/jing_meng_wan_liberal_spring/q20_fig1.png}
\end{problem}

\begin{solution}
\begin{enumerate}
\item
设 \(\Pi\) 为过 \(V\)，\(C\)，\(N\) 的平面. 因为 \(VN\perp\) 平面 \(ABC\)，且 \(AB\perp CD\)，所以直线 \(AB\) 垂直于平面 \(\Pi\). 平面 \(\Pi\) 与平面 \(MAB\) 的交线为 \(DM\)，与平面 \(ABC\) 的交线为 \(DC\)，而这两条交线都过 \(D\)，故 \(\angle MDC\) 正是二面角 \(M-AB-C\) 的平面角.

\item 为证明第二问，建立直角坐标系：取 \(D\) 为原点，\(AB\) 为 \(x\) 轴，平面 \(\Pi\) 为 \(yz\) 平面，并令
\(A=(-u,0,0)\)，\(B=(v,0,0)\)，\(C=(0,c,0)\)，\(N=(0,n,0)\)，\(V=(0,n,w)\)
其中 \(u\)，\(v\)，\(c\)，\(w>0\). 设 \(M=(0,r,s)\). 记 \(\theta=\angle MDC=\angle CVN\). 由 \(\angle MDC=\theta\) 得
\[
\frac{s}{r}=\tan\theta
\]
在直角三角形 \(CVN\) 中，\(\angle CVN=\theta\)，所以
\[
\frac{c-n}{w}=\tan\theta
\]
于是 \(n-c=-w\tan\theta\). 取 \(VC\) 的方向向量 \(\bs{u}=(0,n-c,w)\)，则
\(\overrightarrow{AM}=(u,r,s)\)，\(\overrightarrow{BM}=(-v,r,s)\)
从而
\[
\bs{u}\cdot\overrightarrow{AM}=(n-c)r+ws
=-wr\tan\theta+ws=0
\]
\[
\bs{u}\cdot\overrightarrow{BM}=(n-c)r+ws=0
\]
因此 \(VC\perp AM\)，且 \(VC\perp BM\). 直线 \(AM\)，\(BM\) 相交于 \(M\)，所以 \(VC\perp\) 平面 \(AMB\).

\item 再求第三问. 在平面 \(\Pi\) 内，设 \(K\) 是 \(M\) 在 \(DC\) 上的正投影. 则 \(MK\) 是点 \(M\) 到平面 \(ABC\) 的距离. 由 \(\angle MDC=\theta\)，有
\[
\frac{MK}{DK}=\tan\theta
\]
又因为 \(M\) 在 \(VC\) 上，且直线 \(VC\) 与竖直方向 \(VN\) 的夹角为 \(\theta\)，所以
\[
\frac{MK}{CK}=\cot\theta
\]
设 \(CD=c\)，则 \(DK+CK=DC=c\). 联立上面两式，得
\(DK=c\cos^2\theta\)，\(MK=c\sin\theta\cos\theta\)
由于平面 \(\Pi\perp AB\)，直线 \(AB\) 与 \(VC\) 的距离等于点 \(D\) 到直线 \(VC\) 的距离. 直线 \(VC\) 与 \(DC\) 的夹角为 \(\frac{\pi}{2}-\theta\)，所以
\(h=CD\sin\left(\frac{\pi}{2}-\theta\right)=c\cos\theta\)，\(c=\frac{h}{\cos\theta}\)
四面体 \(MABC\) 的底面取 \(\triangle ABC\)，其面积为 \(\frac12ac\)，高为 \(MK\)，故
\[
V_{MABC}=\frac13\cdot\frac12ac\cdot MK
=\frac{ac^2\sin\theta\cos\theta}{6}
=\frac{ah^2\tan\theta}{6}
\]
\end{enumerate}
\end{solution}

\begin{problem}
某摩托车生产企业，上年度生产摩托车的投入成本为 \(1\text{ 万元/辆}\)，出厂价为 \(1.2\text{ 万元/辆}\)，年销售量为 \(1000\text{ 辆}\). 本年度为适应市场需求，计划提高产品档次，适度增加投入成本. 若每辆车投入成本增加的比例为 \(x\)（\(0 < x < 1\)），则出厂价相应提高的比例为 \(0.75x\). 同时预计年销售量增加的比例为 \(0.6x\). 已知年利润 \(=(\text{出厂价} - \text{投入成本}) \times \text{年销售量}\).

\begin{enumerate}
\item 写出本年度预计的年利润 \(y\) 与投入成本增加的比例 \(x\) 的关系式；
\item 为使本年度的年利润比上年有所增加，问投入成本增加的比例 \(x\) 应在什么范围内？
\end{enumerate}
\end{problem}

\begin{solution}
\begin{enumerate}
\item 每辆车本年度的投入成本为 \(1+x\) 万元，出厂价为
\[
1.2(1+0.75x)=1.2+0.9x
\]
万元，年销售量为 \(1000(1+0.6x)\) 辆. 所以每辆车的利润为
\[
(1.2+0.9x)-(1+x)=0.2-0.1x
\]
万元，从而本年度预计年利润为
\[
y=(0.2-0.1x)\cdot1000(1+0.6x)
=(-100x+200)(1+0.6x)
=-60x^2+20x+200
\]
其中 \(0<x<1\).
\item 上年度年利润为
\[
(1.2-1)\times1000=200
\]
万元. 要使本年度利润增加，需要
\[
-60x^2+20x+200>200
\]
即
\[
20x(1-3x)>0
\]
结合 \(0<x<1\)，得到 \(0<x<\dfrac13\). 因此投入成本增加的比例应满足
\[
0<x<\frac13
\]
\end{enumerate}
\end{solution}

\begin{problem}
已知抛物线 \(y^{2} = 2px\)（\(p > 0\)）. 过动点 \(M(a, 0)\) 且斜率为 1 的直线 \(l\) 与该抛物线交于不同的两点 \(A\)，\(B\).

\begin{enumerate}
\item 若 \(\left|AB\right| \leqslant 2p\)，求 \(a\) 的取值范围；
\item 若线段 \(AB\) 的垂直平分线交 \(AB\) 于点 \(Q\)，交 \(x\) 轴于点 \(N\)，试求 \(\mathrm{Rt}\triangle\mathrm{MNQ}\) 的面积.
\end{enumerate}
\end{problem}

\begin{solution}
直线 \(l\) 过 \(M(a,0)\) 且斜率为 \(1\)，所以方程为
\[
y=x-a,\qquad\text{即}\qquad x=y+a
\]
将其代入抛物线方程，得
\[
y^2=2p(y+a)
\]
即
\[
y^2-2py-2pa=0
\]
设交点 \(A\)，\(B\) 的纵坐标分别为 \(y_1\)，\(y_2\). 因为 \(A\)，\(B\) 是两个不同的交点，所以
\[
\Delta=(-2p)^2-4(-2pa)=4p(p+2a)>0
\]
从而 \(a>-\dfrac p2\). 由根的差的绝对值，
\[
|y_1-y_2|=\sqrt{\Delta}=2\sqrt{p(p+2a)}
\]
直线 \(l\) 的斜率为 \(1\)，故 \(|x_1-x_2|=|y_1-y_2|\)，于是
\[
|AB|=\sqrt{(x_1-x_2)^2+(y_1-y_2)^2}
=\sqrt2\,|y_1-y_2|
=2\sqrt{2p(p+2a)}
\]
若 \(|AB|\leqslant2p\)，则
\[
2\sqrt{2p(p+2a)}\leqslant2p
\]
两边均为非负数，平方并除以 \(4p\) 得
\[
2(p+2a)\leqslant p
\]
即 \(a\leqslant-\dfrac p4\). 与 \(a>-\dfrac p2\) 合并，得到
\[
-\frac p2<a\leqslant-\frac p4
\]
这就是第一问的取值范围.

由韦达定理，\(y_1+y_2=2p\)，所以弦 \(AB\) 的中点为
\[
Q=\left(a+\frac{y_1+y_2}{2},\frac{y_1+y_2}{2}\right)=(a+p,p)
\]
由于 \(AB\) 的斜率为 \(1\)，其垂直平分线斜率为 \(-1\)，且过 \(Q\)，方程为
\[
y-p=-(x-a-p)
\]
即 \(y=-x+a+2p\). 令 \(y=0\)，得到
\[
N=(a+2p,0)
\]
而 \(M=(a,0)\)，所以 \(|MN|=2p\)，点 \(Q\) 到 \(x\) 轴的距离为 \(p\). 因此
\[
S_{\triangle MNQ}=\frac12\cdot2p\cdot p=p^2
\]
故 \(\mathrm{Rt}\triangle MNQ\) 的面积为 \(p^2\).
\end{solution}
